\subsection{对偶性与弱收敛}\label{sec:duality-weak-convergence}

在分布理论中, 已经遇到许多对偶性的思想 (通过研究线性泛函如何作用于函数来研究该函数). 这些思想在 Banach 空间和 Hilbert 空间的研究中也十分有力.

\subsubsection{有界线性映射}\label{subsec:bounded-linear-mappings}

\begin{Definition}
\label{def:bounded-linear-mapping}
设 $X,Y$ 为赋范向量空间. 若存在常数 $C$, 使对每个 $x\in X$ 有 $\|Lx\|\leq C\|x\|$, 则称线性映射 $L:X\to Y$ \textbf{有界}.
\end{Definition}

对线性映射, 连续性和有界性的概念本质上相同.

\begin{Definition}
\label{def:operator-sequential-continuity}
若每当 $x_n\in X$ 为满足 $x_n\to x$ 的序列时都有 $L(x_n)\to L(x)$, 就说算子 $L:X\to Y$ 在点 $x\in X$ 处\textbf{连续}.
\end{Definition}

\begin{Remark}
\label{rem:sequential-versus-topological-continuity}
Definition~\ref{def:operator-sequential-continuity} 常称为\emph{序列连续性}的定义, 以区别于该定义的拓扑版本.
\end{Remark}

\begin{Lemma}
\label{lem:linear-mapping-bounded-continuity-equivalence}
设 $L:X\to Y$ 为线性映射.\begin{enumerate}
\item 若 $L$ 在原点处连续, 则它在每个 $x\in X$ 处连续.
\item $L$ 连续的充分必要条件是它有界.
\end{enumerate}
\end{Lemma}

\begin{Proof}
为证明第 1 部分, 假设 $L$ 在原点处连续. 注意, 算子的线性性意味着 $L(0)=0$. 那么对任意 $x\in X$ 及任意满足 $\lim_{n\to\infty}x_n=x$ 的序列 $x_n\in X$, 有
\begin{equation}\label{eq:linear-map-continuity-at-origin-translation}
\lim_{n\to\infty}L(x_n)-L(x)=\lim_{n\to\infty}L(x_n-x)=L(0)=0,
\end{equation}
因而 $L$ 在 $x$ 处连续.

为证明第 2 部分, 注意有界性的定义意味着在零点处连续. 因而由第 1 部分, 有界性蕴含连续性. 另一方面, 假设 $L$ 不有界. 那么对每个 $n$, 存在 $y_n$, 使 $\|Ly_n\|>n^2\|y_n\|$. 注意 $x_n:=y_n(n\|y_n\|)^{-1}\to0$, 但 $\|Lx_n\|>n$. 因此 $L$ 在原点处不连续. 从而连续性蕴含有界性.
\end{Proof}

自然可以考察所有有界线性映射的集合. 由定义立即可知, 这一集合构成向量空间. 此外, 若取 Definition~\ref{def:bounded-linear-mapping} 中可能的最小常数, 那么这个量就给出线性映射 ``大小'' 的一个度量. 这引出以下定义.

\begin{Definition}
\label{def:bounded-operator-space-norm}
用 $\mathcal{L}(X,Y)$ 表示从 $X$ 到 $Y$ 的所有有界线性映射的集合. 若 $X=Y$, 也将 $\mathcal{L}(X,X)$ 写为 $\mathcal{L}(X)$. 此外, 若 $L\in\mathcal{L}(X,Y)$, 令
\begin{equation}\label{eq:bounded-linear-map-operator-norm}
\|L\|:=\sup_{\|x\|=1}\dfrac{\|Lx\|}{\|x\|}=\sup_{\|x\|\ne0}\dfrac{\|Lx\|}{\|x\|}.
\end{equation}

\end{Definition}

\begin{Theorem}
\label{thm:bounded-operator-space-completeness}
设 $X,Y$ 为 Banach 空间. 那么 $\mathcal{L}(X,Y)$ 在 \eqref{eq:bounded-linear-map-operator-norm} 定义的范数下也是 Banach 空间.
\end{Theorem}

证明是直接的, 留作习题 (Problem~\ref{ex:bounded-operator-space-completeness}). 线性映射也称线性算子, 将在下一章中作更深入研究. 本节关注从 Banach 空间到其标量域的线性映射这一特殊情形.

\begin{Definition}
\label{def:linear-functional-dual-space}
设 $X$ 为实 (复) Banach 空间. 那么 $X$ 上的\textbf{线性泛函}是从 $X$ 到 $\mathbb{R}$ ($\mathbb{C}$) 的有界线性映射. $X$ 上所有线性泛函组成的空间称为 $X$ 的\textbf{对偶空间}, 记为 $X^*$.
\end{Definition}
